\section{A real ray host with controlled ramification orbits}
\label{sec:cubic-genus-host-research}

This is a compatible host construction for a proposed genus-class
argument.  The required isotypic quotient of the \(p\)-split class group
is not proved in this note.  In particular, this note alone supplies
neither strict class positivity nor the numerical frame inequalities.

\begin{proposition}[AC-G92-host: a cyclic real host with prescribed split ramification]
\label{AC-G92-host}
Let \(p\) be odd.  Let \(F/\mathbb Q\) be a totally real Galois field
and let \(M/\mathbb Q\) be finite Galois, with \(F\subset M\) and
\(p\nmid[M:F]\).  Let \(S_F\) be the primes of \(F\) above \(p\), and
\(S_M\) the primes of \(M\) above them.  Write
\[
 E_{A,S}=\mathcal O_A[1/p]^\times,\qquad
 h_{F,S}(p)=\#\operatorname{Cl}(\mathcal O_F[1/p])(p).
\]
For any \(r\geq1\) there are distinct fresh rational primes \(q_i\),
places \(v_i\) of \(F\) above \(q_i\), and a totally real cyclic
extension \(T'/F\) with the following properties:
\[
 [T':F]=p^s\leq p\,h_{F,S}(p),\qquad s\geq1.
 \tag{G92.1}
\]
Every \(q_i\) splits completely in \(M\).  The extension \(T'/F\)
is ramified only at the selected \(v_i\), with inertia exactly of
order \(p\) at each, and every place in \(S_F\) splits completely.
The fields \(T'\) and \(M\) are linearly disjoint over \(F\).
For \(M'=MT'\), every element of \(E_{M,S_M}\) is a local norm
from \(M'\) at every place of \(M\).

With \(\Delta=\operatorname{Gal}(M/F)\), the inertia module at these
ramified places is exactly
\[
 \bigoplus_{i=1}^r\ \bigoplus_{w\mid v_i}I_w(M'/M)
       \simeq \mathbb F_p[\Delta]^{\,r}
 \tag{G92.2}
\]
as a \(\Delta\)-module.  The degree bound in (G92.1) is independent
of \(r\).  No conclusion about the image of (G92.2) in a class
group is included in this proposition.
\end{proposition}

\begin{proof}
\emph{Choose the primes before the ray character.}
The finitely generated group \(E_{M,S_M}\) has a finite set of
generators including its torsion subgroup.  Take a finite Galois
extension \(\mathcal K/\mathbb Q\) containing \(M,\mu_p\), and the
\(p\)-th roots of these generators.  Enlarge it to contain the
\(p\)-th roots of a finite generating set of \(E_{F,S_F}\), or note
that these are already \(S_M\)-units.  A Galois closure is harmless.
Chebotarev supplies arbitrarily many rational primes splitting
completely in \(\mathcal K\), outside any prescribed finite set.
Choose \(q_i\) there and one \(v_i\mid q_i\) in \(F\).
Then \(\mathcal O_F/v_i=\mathbb F_{q_i}\), \(p\mid(q_i-1)\), and
every \(F\)- or \(M\)-\(S\)-unit has \(p\)-th-power reduction at
every place above \(v_i\).  Roots of torsion units are included;
no assertion that \(v_p(q_i-1)=1\) is made.

\emph{Use the \(S\)-ray class group.}
Put \(\mathfrak m=\prod_i v_i\), with no infinite modulus, and
let \(\operatorname{Cl}_{\mathfrak m,S}(F)\) be the ray class group
modulo the additional ideal classes of all primes in \(S_F\).
The ordinary ideal definition of ray classes gives the exact sequence
\[
 E_{F,S_F}\longrightarrow
 \prod_i(\mathcal O_F/v_i)^\times
 \longrightarrow \operatorname{Cl}_{\mathfrak m,S}(F)
 \longrightarrow \operatorname{Cl}(\mathcal O_F[1/p])
 \longrightarrow0.
 \tag{G92.3}
\]
For clarity, this is the usual residue-unit/ray-class sequence with
the ideals at \(S_F\) made trivial.  If a product of such ideals
becomes principal, its generator is precisely an \(S_F\)-unit;
this accounts for the first term and proves exactness at the
residue-unit term without treating ordinary units as \(S\)-units.

Choose a nontrivial character
\(\theta_i:(\mathcal O_F/v_i)^\times\to\mu_p\) for each \(i\).
Their product is trivial on the first term of (G92.3), because every
such residue is a \(p\)-th power.  It therefore gives a character
of the image \(C\) of the residue-unit group, nontrivial of order
\(p\) on each local factor and with total image of order \(p\).
Pass to the \(p\)-primary part
\(A=\operatorname{Cl}_{\mathfrak m,S}(F)(p)\).
Extend this character of \(C(p)\) to \(A\) with values in
\(\mathbb Q_p/\mathbb Z_p\).  Such an extension exists because this
target is divisible; it can also be constructed by successively
adjoining one generator of a finite abelian group and choosing a
root of its prescribed multiple.
Its finite image is cyclic of order \(p^s\), and
\[
 p^{s-1}=\#(\operatorname{im}\theta/\theta(C))
           \leq \#(A/C(p))=h_{F,S}(p).
 \tag{G92.4}
\]
This does not assert an injection of a particular order-\(p\)
subgroup of the residue group: torsion \(S\)-units may meet that
subgroup.  It is the order-\(p\) residue \emph{character}, trivial
on all \(S\)-units, which is extended.

Class field theory produces the cyclic extension \(T'/F\).
The finite modulus and the local restrictions show that it is
unramified outside the \(v_i\), with inertia exactly \(p\) at each.
Quotienting by the \(S_F\)-ideal classes forces their Frobenius
elements to be trivial, so these places split completely.
There is no ramification at infinity in an odd-degree extension
of a totally real field.  Thus \(T'\) is totally real.
Since \([M:F]\) is prime to \(p\), linear disjointness is automatic.

\emph{Local norms and the ramification module.}
At a place \(w\mid v_i\), the completion \(M_w\) is \(F_{v_i}\).
The local norm subgroup on units is the kernel of the order-\(p\)
character \(\theta_i\) on the residue field; principal units are in
the kernel.  The reductions of all \(E_{M,S_M}\) elements are
\(p\)-th powers, so those elements are local norms.
At every other finite place not in \(S_M\), the local extension is
unramified and a global \(S_M\)-unit is a unit; the norm on units
of an unramified local extension is surjective.
At the places in \(S_M\) the extension splits, and the infinite
places impose no condition here.  This proves the asserted
local-norm statement.  A use of the cyclic Hasse norm theorem could
upgrade it to a global element-norm statement, but is unnecessary
for the proposition.

Finally
\(\operatorname{Gal}(M'/F)=\Delta\times\operatorname{Gal}(T'/F)\).
For each \(v_i\), its primes in \(M\) form one free transitive
\(\Delta\)-set.  Their inertia groups in \(M'/M\) are all the
unique order-\(p\) subgroup of the cyclic second factor.
Conjugation by \(\Delta\) acts trivially on that factor and
permutes the places.  This proves (G92.2).
\end{proof}

\paragraph{Compatibility with the actual cubic automorphic host.}
Take \(F=H^+\) for the actual G83/G87 host, and enlarge its
prime-to-\(p\) residual splitting extension \(M/F\) as above.
Put \(H'=HT'\).  Then \(H'/T'\) is the base change of \(H/F\);
it is CM, unramified at every finite place, and relatively split
at the coefficient-prime and old automorphic ramification places.
Every coefficient-prime completion is unchanged, since \(S_F\)
splits in \(T'\).
The extension \(H'/H\) is cyclic of \(p\)-power degree.
The existing absolutely irreducible rank-three representation
remains absolutely irreducible by the determinant/self-twist
argument for successive degree-\(p\) restrictions, and the already
retained cyclic automorphic base-change theorem gives its
cuspidal realization.  Ordinary flags and exact finite away
types restrict.  The new primes \(v_i\) create no representation
ramification, because the old representation was unramified
there.

Although \(T'/\mathbb Q\) need not be Galois, its normal closure
over the old Galois field \(F\) is a compositum of cyclic
\(p\)-power extensions of \(F\), and therefore has \(p\)-power
degree over \(F\).  The normal closure of \(H'\) is accordingly
a \(p\)-extension of \(H\).  Its intersection with the
prime-to-\(p\) field \(\Omega\) used in G89 is still \(K\):
every map to a prime-to-\(p\) Galois quotient kills that normal
\(p\)-subgroup.  The same scalar residual auxiliary class is
therefore compatible with complete splitting in this new normal
closure.  A fresh auxiliary place and the retained native
propositions give support just as in G89.

Also \([H':\mathbb Q]=[H:\mathbb Q]p^s\), whose only prime
factors are \(2\) and \(269\); thus \(H'\) cannot contain
\(\mu_{269}\), whose degree is \(268=4\cdot67\).
At the coefficient prime the local fields are literally unchanged,
so their torsion and full local character interfaces are unchanged
as well.
For the even rational dual used in G91, base change of
\(T'/F\) after adjoining its abelian splitting field and \(\mu_p\)
is again a \(p\)-extension.  The already-audited finite-\(p\)
\(\mu=0\) ascent therefore continues to apply.

\paragraph{Why a selected place matters.}
One could instead take a rational cyclic degree-\(p\) field ramified
at rational primes \(q_i\equiv1\pmod p\), using their order-\(p\)
Dirichlet characters.  That also produces an admissible real host,
but it ramifies \emph{all} the places of \(F\) above every \(q_i\).
The resulting permutation module has \([F:\mathbb Q]\) copies
per rational prime.  A strong genus quotient of that size would
be unsuitable for the source's small-class-sector frame inequality.
Construction (G92.3) selects one place of \(F\) at a time and gives
exactly one copy in (G92.2).

\paragraph{Retained inputs and unresolved arithmetic step.}
The inputs are the finite generation of \(S\)-units, qualitative
Chebotarev, the ordinary ray class exact sequence and class field
existence, and local reciprocity with its unit-norm interpretation.
The localization in (G92.3), character extension and all host
comparisons have been proved above.  The ray-class and local
reciprocity inputs retain the previously audited Milne class-field
scope, rather than claiming a new proof of global reciprocity.
For this draft the cached older primary mirror was read afresh:
Milne, \emph{Class Field Theory}, V.1.5 and its entire printed proof
(pp.~115--116), and the statements of V.3.5--3.8 (pp.~123--124).
The \(S\)-localization in (G92.3) is the new elementary deduction
from that sequence.  The earlier local-reciprocity scope is retained.
No bound for the chosen \(q_i\) is asserted.
An actual \(\Delta\)-isotypic quotient of
\(\operatorname{Cl}(\mathcal O_{M'}[1/p])/p\), including the
unit and \(p\)-splitting corrections, is still required.
Neither positivity nor a frame inequality is inferred from the
ramification permutation module alone.

\paragraph{Geometric dictionary.}
The ray extensions give finite covers of rings of integers, finite
\'{e}tale after removing the chosen tame ramification places.
Making the coefficient-prime Frobenius classes trivial gives complete
splitting over those points, while (G92.2) records the permutation of
the tame inertia fibers under the residual cover.  This uses the
already-audited finite-cover and fundamental-group dictionary; it
does not assert a new archival SGA reading or identify inertia fibers
with class-group sectors before the separate genus argument.
